% GENERATED FILE. Edit canonical lessons, then run npm run generate:books.
% canonical-source-hash: fnv1a64:a6e9e3c9976cc623
% canonical-lessons: GU-C487-cakri, GU-C487-dilaso, GU-C487-malis, GU-C487-dabvu, GU-C487-odhadvu

\chapter{Caring for Grandmother}
\label{ch:gu-caring-for-grandmother}
\hlchaptermodality{487}

\begin{chapteropening}
\textbf{By the end of this chapter:} \emph{I can say how we look after a sick relative.}

Complete the last lesson of chapter 487: I can say how we look after a sick relative.
\end{chapteropening}

\section[\texorpdfstring{\gu{ચાકરી} (cākrī)}{ચાકરી (cākrī)}]{\gu{ચાકરી} (cākrī) — nursing care; looking after}

\label{lesson:GU-C487-cakri}



\begin{quote}
Before the new one: say the Gujarati for diet, then the Gujarati for to waste away.
\end{quote}

\subsection*{You'll want to know: \gu{ચાકરી}}
\textbf{\gu{ચાકરી}} — \emph{cākrī} — \textquotedblleft{}nursing care; looking after\textquotedblright{}.

\begin{quote}
\gu{વહુએ} \gu{સાસુની} \gu{ઘણી} \gu{ચાકરી} \gu{કરી}.
\end{quote}

\begin{quote}
\emph{vahue sāsunī ghaṇī cākrī karī.}
\end{quote}

\begin{quote}
The daughter-in-law looked after her mother-in-law devotedly.
\end{quote}

A feminine noun.

\begin{cousinweb}
\textbf{\gu{ચાકરી}} (\emph{cākrī}) is usually traced to Persian \emph{čākarī}, \textquotedblleft{}service\textquotedblright{}, built on \emph{čākar}, \textquotedblleft{}a servant\textquotedblright{}.
\end{cousinweb}

\begin{grammarlens}[title={What you've built}]
One more word for describing what you lived and felt.
\end{grammarlens}

\subsection*{Your turn}
\begin{itemize}
\raggedright
  \item \emph{Say it:} \emph{cākrī}
  \item \emph{Say it:} \emph{cākrī}, once more
  \item \emph{Say it:} \emph{vahue sāsunī ghaṇī cākrī karī.} — The daughter-in-law looked after her mother-in-law devotedly.
  \item \emph{Recall:} say the Gujarati for diet, then the Gujarati for to waste away, then say \emph{cākrī} again
\end{itemize}

\begin{tcolorbox}[breakable,colback=teal!4,colframe=teal!35!black,title={Before you move on}]
What does \gu{ચાકરી} mean? (\textquotedblleft{}Nursing care; looking after\textquotedblright{}.) Say it once more.
\end{tcolorbox}
\section[\texorpdfstring{\gu{દિલાસો} (dilāso)}{દિલાસો (dilāso)}]{\gu{દિલાસો} (dilāso) — comfort; consoling words}

\label{lesson:GU-C487-dilaso}



\begin{quote}
Before the new one: say the Gujarati for to waste away, then the Gujarati for nursing care.
\end{quote}

\subsection*{You'll want to know: \gu{દિલાસો}}
\textbf{\gu{દિલાસો}} — \emph{dilāso} — \textquotedblleft{}comfort; consoling words\textquotedblright{}.

\begin{quote}
\gu{મિત્રે} \gu{મને} \gu{દિલાસો} \gu{આપ્યો}.
\end{quote}

\begin{quote}
\emph{mitre mane dilāso āpyo.}
\end{quote}

\begin{quote}
My friend comforted me.
\end{quote}

A masculine noun.

\begin{cousinweb}
\textbf{\gu{દિલાસો}} (\emph{dilāso}) is usually traced to Persian \emph{dilāsā}, \textquotedblleft{}consolation\textquotedblright{}, built on \emph{dil}, \textquotedblleft{}heart\textquotedblright{}.
\end{cousinweb}

\begin{grammarlens}[title={What you've built}]
One more word for describing what you lived and felt.
\end{grammarlens}

\subsection*{Your turn}
\begin{itemize}
\raggedright
  \item \emph{Say it:} \emph{dilāso}
  \item \emph{Say it:} \emph{dilāso}, once more
  \item \emph{Say it:} \emph{mitre mane dilāso āpyo.} — My friend comforted me.
  \item \emph{Recall:} say the Gujarati for to waste away, then the Gujarati for nursing care, then say \emph{dilāso} again
\end{itemize}

\begin{tcolorbox}[breakable,colback=teal!4,colframe=teal!35!black,title={Before you move on}]
What does \gu{દિલાસો} mean? (\textquotedblleft{}Comfort; consoling words\textquotedblright{}.) Say it once more.
\end{tcolorbox}
\section[\texorpdfstring{\gu{માલિશ} (māliś)}{માલિશ (māliś)}]{\gu{માલિશ} (māliś) — a massage}

\label{lesson:GU-C487-malis}



\begin{quote}
Before the new one: say the Gujarati for nursing care, then the Gujarati for consoling words.
\end{quote}

\subsection*{You'll want to know: \gu{માલિશ}}
\textbf{\gu{માલિશ}} — \emph{māliś} — \textquotedblleft{}a massage\textquotedblright{}.

\begin{quote}
\gu{દાદી} \gu{બાળકને} \gu{રોજ} \gu{તેલથી} \gu{માલિશ} \gu{કરે} \gu{છે}.
\end{quote}

\begin{quote}
\emph{dādī bāḷakne roj telthī māliś kare che.}
\end{quote}

\begin{quote}
Grandmother massages the baby with oil every day.
\end{quote}

A feminine noun.

\begin{cousinweb}
\textbf{\gu{માલિશ}} (\emph{māliś}) is usually traced to Persian \emph{māliš}, \textquotedblleft{}rubbing\textquotedblright{}.
\end{cousinweb}

\begin{grammarlens}[title={What you've built}]
One more word for describing what you lived and felt.
\end{grammarlens}

\subsection*{Your turn}
\begin{itemize}
\raggedright
  \item \emph{Say it:} \emph{māliś}
  \item \emph{Say it:} \emph{māliś}, once more
  \item \emph{Say it:} \emph{dādī bāḷakne roj telthī māliś kare che.} — Grandmother massages the baby with oil every day.
  \item \emph{Recall:} say the Gujarati for nursing care, then the Gujarati for consoling words, then say \emph{māliś} again
\end{itemize}

\begin{tcolorbox}[breakable,colback=teal!4,colframe=teal!35!black,title={Before you move on}]
What does \gu{માલિશ} mean? (\textquotedblleft{}A massage\textquotedblright{}.) Say it once more.
\end{tcolorbox}
\section[\texorpdfstring{\gu{દાબવું} (dābvũ)}{દાબવું (dābvũ)}]{\gu{દાબવું} (dābvũ) — to massage by pressing; to press (the legs) to ease pain}

\label{lesson:GU-C487-dabvu}



\begin{quote}
Before the new one: say the Gujarati for consoling words, then the Gujarati for a massage.
\end{quote}

\subsection*{You'll want to know: \gu{દાબવું}}
\textbf{\gu{દાબવું}} — \emph{dābvũ} — \textquotedblleft{}to massage by pressing; to press (the legs) to ease pain\textquotedblright{}.

\begin{quote}
\gu{રાતે} \gu{મીરા} \gu{દાદીના} \gu{પગ} \gu{દાબે} \gu{છે}.
\end{quote}

\begin{quote}
\emph{rāte Mīrā dādīnā pag dābe che.}
\end{quote}

\begin{quote}
At night Mira presses her grandmother's legs to ease them.
\end{quote}

Past, with \gu{તેણે} (he or she …): \gu{તેણે} \gu{દાદીના} \gu{પગ} \gu{દાબ્યા}, \emph{teṇe dādīnā pag dābyā}.

\begin{cousinweb}
\textbf{\gu{દાબવું}} (\emph{dābvũ}) goes back to a reconstructed \emph{*dabb-}, \textquotedblleft{}press\textquotedblright{}, which some trace to Dravidian; CDIAL lists the Gujarati verb, \textquotedblleft{}to press\textquotedblright{} (CDIAL 6173).
\end{cousinweb}

\begin{grammarlens}[title={What you've built}]
One more word for describing what you lived and felt.
\end{grammarlens}

\subsection*{Your turn}
\begin{itemize}
\raggedright
  \item \emph{Say it:} \emph{dābvũ}
  \item \emph{Say it:} \emph{dābvũ}, once more
  \item \emph{Say it:} \emph{rāte Mīrā dādīnā pag dābe che.} — At night Mira presses her grandmother's legs to ease them.
  \item \emph{Recall:} say the Gujarati for consoling words, then the Gujarati for a massage, then say \emph{dābvũ} again
\end{itemize}

\begin{tcolorbox}[breakable,colback=teal!4,colframe=teal!35!black,title={Before you move on}]
What does \gu{દાબવું} mean? (\textquotedblleft{}To massage by pressing; to press (the legs) to ease pain\textquotedblright{}.) Say it once more.
\end{tcolorbox}
\section[\texorpdfstring{\gu{ઓઢાડવું} (oḍhāḍvũ)}{ઓઢાડવું (oḍhāḍvũ)}]{\gu{ઓઢાડવું} (oḍhāḍvũ) — to cover someone with a blanket}

\label{lesson:GU-C487-odhadvu}



\begin{quote}
Before the new one: say the Gujarati for a massage, then the Gujarati for to massage by pressing.
\end{quote}

\subsection*{You'll want to know: \gu{ઓઢાડવું}}
\textbf{\gu{ઓઢાડવું}} — \emph{oḍhāḍvũ} — \textquotedblleft{}to cover someone with a blanket\textquotedblright{}.

\begin{quote}
\gu{માતાએ} \gu{બાળકને} \gu{ધાબળો} \gu{ઓઢાડ્યો}.
\end{quote}

\begin{quote}
\emph{mātāe bāḷakne dhābaḷo oḍhāḍyo.}
\end{quote}

\begin{quote}
The mother tucked a blanket around the child.
\end{quote}

Past, with \gu{તેણે} (he or she …): \gu{તેણે} \gu{દાદીને} \gu{ધાબળો} \gu{ઓઢાડ્યો}, \emph{teṇe dādīne dhābaḷo oḍhāḍyo}.

\begin{cousinweb}
\textbf{\gu{ઓઢાડવું}} (\emph{oḍhāḍvũ}) is the causative of \textbf{\gu{ઓઢવું}} (\emph{oḍhvũ}), \textquotedblleft{}to wrap oneself in\textquotedblright{}, which goes back to a reconstructed \emph{*ōḍḍh-}, \textquotedblleft{}put on, cover\textquotedblright{} (CDIAL 2547).
\end{cousinweb}

\begin{grammarlens}[title={What you've built}]
One more word for describing what you lived and felt.
\end{grammarlens}

\subsection*{Your turn}
\begin{itemize}
\raggedright
  \item \emph{Say it:} \emph{oḍhāḍvũ}
  \item \emph{Say it:} \emph{oḍhāḍvũ}, once more
  \item \emph{Say it:} \emph{mātāe bāḷakne dhābaḷo oḍhāḍyo.} — The mother tucked a blanket around the child.
  \item \emph{Recall:} say the Gujarati for a massage, then the Gujarati for to massage by pressing, then say \emph{oḍhāḍvũ} again
\end{itemize}

\begin{tcolorbox}[breakable,colback=teal!4,colframe=teal!35!black,title={Before you move on}]
What does \gu{ઓઢાડવું} mean? (\textquotedblleft{}To cover someone with a blanket\textquotedblright{}.) Say it once more.
\end{tcolorbox}
